\chapter{Background and Related Work}
\label{ch:background}

\section{Petri Net Fundamentals}
\label{sec:pn-fundamentals}

\subsection{Classical Petri Nets}
\label{subsec:classical-pn}

\textbf{Petri nets}, introduced by Carl Adam Petri in 1962 \citep{Petri1962}, are a mathematical formalism for modeling concurrent systems. A \textbf{classical Petri net} is a 5-tuple:

\begin{equation}
PN = (P, T, F, W, M_0)
\end{equation}

Where:

\begin{itemize}
    \item $P$: Finite set of \textbf{places} (circles, represent conditions or resources)
    \item $T$: Finite set of \textbf{transitions} (rectangles, represent events or actions)
    \item $F \subseteq (P \times T) \cup (T \times P)$: \textbf{Flow relation} (arcs connecting places to transitions and vice versa)
    \item $W: F \to \mathbb{N}^+$: \textbf{Weight function} (arc multiplicities, stoichiometric coefficients)
    \item $M_0: P \to \mathbb{N}$: \textbf{Initial marking} (token distribution at time $t=0$)
\end{itemize}

\textbf{Graphical notation}:

\begin{itemize}
    \item \textbf{Places}: Circles containing tokens (black dots)
    \item \textbf{Transitions}: Rectangles or bars
    \item \textbf{Arcs}: Arrows with weights (omitted if weight = 1)
\end{itemize}

\textbf{Example} (Simple chemical reaction \ce{A + B -> C}):

\begin{itemize}
    \item Places: $P = \{A, B, C\}$
    \item Transitions: $T = \{T_1\}$
    \item Arcs: $F = \{(A,T_1), (B,T_1), (T_1,C)\}$
    \item Weights: $W(A,T_1)=1$, $W(B,T_1)=1$, $W(T_1,C)=1$
    \item Initial marking: $M_0(A)=2$, $M_0(B)=1$, $M_0(C)=0$
\end{itemize}

\subsection{Petri Net Semantics}
\label{subsec:pn-semantics}

\textbf{Enabling}: Transition $t$ is \textbf{enabled} at marking $M$ if:

\begin{equation}
\forall p \in \bullet t: M(p) \geq W(p,t)
\end{equation}

Where $\bullet t$ denotes the \textbf{preset} (input places) of $t$.

\textbf{Firing}: When enabled transition $t$ fires:

\begin{enumerate}
    \item \textbf{Consume tokens} from input places: $M'(p) = M(p) - W(p,t)$ for all $p \in \bullet t$
    \item \textbf{Produce tokens} at output places: $M'(p) = M(p) + W(t,p)$ for all $p \in t\bullet$
\end{enumerate}

Where $t\bullet$ denotes the \textbf{postset} (output places) of $t$.

\textbf{Example} (\ce{A + B -> C}):

\begin{itemize}
    \item \textbf{Before firing}: $M = \{A:2, B:1, C:0\}$
    \item \textbf{$T_1$ enabled?} Yes ($M(A)=2 \geq 1$, $M(B)=1 \geq 1$)
    \item \textbf{$T_1$ fires}: Consume 1 token from A, 1 from B, produce 1 at C
    \item \textbf{After firing}: $M' = \{A:1, B:0, C:1\}$
\end{itemize}

\subsection{Petri Net Properties}
\label{subsec:pn-properties}

\textbf{Behavioral properties}:

\begin{itemize}
    \item \textbf{Reachability}: Can marking $M'$ be reached from $M_0$?
    \item \textbf{Boundedness}: Is there a maximum token count per place?
    \item \textbf{Liveness}: Can all transitions eventually fire?
    \item \textbf{Deadlock-freedom}: Does the net avoid states where no transition is enabled?
\end{itemize}

\textbf{Structural properties} (independent of $M_0$):

\begin{itemize}
    \item \textbf{P-invariants} (place invariants): Token-conserving sets of places
    \begin{itemize}
        \item Linear combination: $y^T \cdot M = y^T \cdot M_0$ (constant sum)
        \item Example: Enzyme conservation ($E_{\text{free}} + E_{\text{bound}} = \text{constant}$)
    \end{itemize}
    \item \textbf{T-invariants} (transition invariants): Firing sequences returning to same marking
    \begin{itemize}
        \item Example: Catalytic cycles (net effect = 0)
    \end{itemize}
\end{itemize}

\textbf{Complexity}:

\begin{itemize}
    \item \textbf{Reachability}: EXPSPACE-complete (decidable but computationally hard)
    \item \textbf{Boundedness}: EXPSPACE-complete
    \item \textbf{Liveness}: EXPSPACE-complete
\end{itemize}

\subsection{Limitations for Biological Modeling}
\label{subsec:pn-limitations}

\textbf{Classical Petri nets have fundamental limitations}:

\begin{enumerate}
    \item \textbf{No regulatory arcs}: Cannot represent catalysis (enzyme unchanged) or inhibition (threshold blocking) without encoding in external code
    \item \textbf{Homogeneous transitions}: All transitions fire instantaneously (no continuous kinetics) or all are timed (no stochastic bursts)
    \item \textbf{Strong independence only}: Parallelism requires disjoint neighborhoods (no shared places)
    \item \textbf{Token abstraction}: Tokens are abstract units (not molecules with formulas)
    \item \textbf{No rate semantics}: Firing frequency is non-deterministic or requires external rate assignment
\end{enumerate}

\textbf{These gaps motivate extensions} for systems biology.

\section{Petri Net Extensions}
\label{sec:pn-extensions}

\subsection{Colored Petri Nets (CPN)}
\label{subsec:colored-pn}

\citet{Jensen1981} introduced \textbf{colored Petri nets} where tokens have \textbf{data types} (colors):

\begin{itemize}
    \item Places hold \textbf{multisets of colored tokens} (e.g., \{glucose:5, fructose:2\})
    \item Transitions have \textbf{guards} (boolean predicates) and \textbf{arc expressions} (token transformations)
    \item Enables modeling complex data structures (proteins with phosphorylation states)
\end{itemize}

\textbf{Advantages}:

\begin{itemize}
    \item \textbf{Compactness}: One place can represent multiple species (different colors)
    \item \textbf{Data manipulation}: Arc expressions compute token transformations
\end{itemize}

\textbf{Limitations for biology}:

\begin{itemize}
    \item \textbf{Complexity}: Guards and expressions are code (not visual arcs)
    \item \textbf{Analysis difficulty}: Structural properties (P-invariants) are harder to compute
    \item \textbf{No heterogeneous dynamics}: Still lacks native continuous/stochastic/burst transitions
\end{itemize}

\subsection{Timed Petri Nets}
\label{subsec:timed-pn}

\citet{Ramchandani1974} added \textbf{time delays} to transitions:

\begin{itemize}
    \item Transitions fire after a \textbf{deterministic delay} (e.g., 5 seconds)
    \item Models sequential processes with known durations
\end{itemize}

\citet{Merlin1974} introduced \textbf{time intervals}:

\begin{itemize}
    \item Transitions fire within a time window $[t_{\min}, t_{\max}]$
    \item Non-deterministic timing
\end{itemize}

\textbf{Limitations}:

\begin{itemize}
    \item \textbf{Still discrete}: No continuous state evolution (ODEs)
    \item \textbf{Homogeneous}: All transitions timed (cannot mix with stochastic)
\end{itemize}

\subsection{Stochastic Petri Nets (SPN)}
\label{subsec:stochastic-pn}

\citet{Molloy1981} introduced \textbf{stochastic Petri nets}:

\begin{itemize}
    \item Transitions fire after \textbf{exponentially distributed delays} (rate $\lambda$)
    \item Equivalent to continuous-time Markov chains (CTMC)
    \item Enables performance analysis (throughput, mean response time)
\end{itemize}

\textbf{Generalized Stochastic Petri Nets (GSPN)} \citep{Ajmone1984}:

\begin{itemize}
    \item \textbf{Immediate transitions} (fire instantly, priority 0)
    \item \textbf{Timed transitions} (exponential delays)
\end{itemize}

\textbf{Applications to biology}:

\begin{itemize}
    \item \textbf{Gene expression}: Transcription/translation as stochastic events
    \item \textbf{Molecular interactions}: Binding/unbinding with rate constants
\end{itemize}

\textbf{Limitations}:

\begin{itemize}
    \item \textbf{Exponential assumption}: Not all biological processes are memoryless
    \item \textbf{State space explosion}: Large models ($>100$ places) are intractable
    \item \textbf{No continuous kinetics}: Cannot model enzyme saturation (Michaelis-Menten)
\end{itemize}

\subsection{Hybrid Petri Nets}
\label{subsec:hybrid-pn}

\citet{David1992} introduced \textbf{continuous and hybrid Petri nets}:

\begin{itemize}
    \item \textbf{Continuous places}: Hold real-valued token counts (not integers)
    \item \textbf{Continuous transitions}: Fire continuously with rate $dM/dt = v(M)$
    \item \textbf{Hybrid nets}: Mix discrete and continuous places/transitions
\end{itemize}

\textbf{Semantics}:

\begin{equation}
\frac{dM(p)}{dt} = \sum_{t} v_t \cdot W(t,p) - \sum_{t} v_t \cdot W(p,t)
\end{equation}

where rate function $v_t$ depends on marking (e.g., mass action: $v = k \cdot M(\text{substrate})$).

\textbf{Advantages for biology}:

\begin{itemize}
    \item \textbf{Enzyme kinetics}: Continuous transitions model Michaelis-Menten rates
    \item \textbf{Multi-scale}: Proteins (discrete) + metabolites (continuous) in one model
\end{itemize}

\textbf{Limitations}:

\begin{itemize}
    \item \textbf{No weak independence theory}: Parallel execution with shared places undefined
    \item \textbf{Limited arc types}: No explicit test/inhibitor arc semantics
    \item \textbf{No biochemical formulas}: Tokens are still abstract (no elemental composition)
\end{itemize}

\subsection{Comparison Summary}
\label{subsec:pn-comparison}

Table~\ref{tab:pn-extensions} summarizes key features of Petri net extensions.

\begin{table}[htbp]
\centering
\caption{Comparison of Petri net extensions for biological modeling.}
\label{tab:pn-extensions}
\begin{tabular}{lccccc}
\toprule
\textbf{Feature} & \textbf{Classical} & \textbf{Colored} & \textbf{Timed} & \textbf{Stochastic} & \textbf{Hybrid} \\
\midrule
Token types & Abstract & Data types & Abstract & Abstract & Real-valued \\
Dynamics & Discrete & Discrete & Timed & Stochastic & Continuous + Discrete \\
Regulation & No & Guards (code) & No & No & Rate functions \\
Parallelism & Strong indep. & Strong indep. & Strong indep. & Strong indep. & Undefined \\
Biological fit & Low & Medium & Medium & Medium & High \\
\bottomrule
\end{tabular}
\end{table}

\textbf{Gap}: No existing extension provides \textbf{weak independence}, \textbf{arc-level regulation}, \textbf{biochemical formulas}, and \textbf{heterogeneous transition coexistence}.

\section{Biological Petri Nets (Bio-PN)}
\label{sec:bio-pn}

\subsection{Early Work (1990s)}
\label{subsec:early-bio-pn}

\textbf{\citet{Reddy1993}}: ``Petri Net Representations in Metabolic Pathways''

\begin{itemize}
    \item First formal application of Petri nets to biochemistry
    \item Modeled glycolysis (10 reactions) as classical Petri net
    \item \textbf{Contribution}: Showed Petri nets can represent metabolic stoichiometry
    \item \textbf{Limitation}: No kinetics (only qualitative), no regulation
\end{itemize}

\textbf{\citet{Hofestaedt1994}}: ``A Petri Net Application to Model Metabolic Processes''

\begin{itemize}
    \item Extended Reddy's work with enzyme catalysis
    \item Used \textbf{colored tokens} to distinguish substrates/products
    \item \textbf{Limitation}: Still no quantitative kinetics
\end{itemize}

\subsection{Qualitative Bio-PNs (2000s)}
\label{subsec:qualitative-bio-pn}

\textbf{\citet{Heiner2008}}: ``Petri Nets for Systems and Synthetic Biology''

\begin{itemize}
    \item Developed \textbf{qualitative Bio-PNs} for signaling networks
    \item Encoded regulatory logic (activation, inhibition) as arc types
    \item \textbf{Tool}: Snoopy software (still actively maintained)
    \item \textbf{Contribution}: Standardized notation for biological arcs
    \item \textbf{Limitation}: Qualitative only (no kinetic parameters)
\end{itemize}

\textbf{\citet{Koch2011}}: ``Modeling Threshold Effects in Biological Systems''

\begin{itemize}
    \item Introduced \textbf{read arcs} (test arcs) and \textbf{inhibitor arcs}
    \item Threshold values for inhibition (e.g., $M(p) < 5$ blocks transition)
    \item \textbf{Limitation}: Thresholds are constants (not Hill equations), no heterogeneous transitions
\end{itemize}

\subsection{Quantitative Bio-PNs}
\label{subsec:quantitative-bio-pn}

\textbf{\citet{Matsuno2003}}: ``Hybrid Petri Net Representation of Gene Regulatory Networks''

\begin{itemize}
    \item Hybrid Petri nets for gene regulation
    \item Continuous places (protein concentrations) + discrete places (gene states)
    \item \textbf{Tool}: Cell Illustrator (commercial)
    \item \textbf{Contribution}: First hybrid Bio-PN with kinetics
    \item \textbf{Limitation}: Proprietary formalism, no weak independence theory
\end{itemize}

\textbf{\citet{Gilbert2006}}: ``From Petri Nets to Differential Equations''

\begin{itemize}
    \item Automatic translation of continuous Petri nets to ODE systems
    \item Mass action kinetics assumed
    \item \textbf{Contribution}: Formalized PN$\to$ODE mapping
    \item \textbf{Limitation}: Only mass action (no Michaelis-Menten, Hill), no stochastic bursts
\end{itemize}

\subsection{Recent Developments (2010s-2020s)}
\label{subsec:recent-bio-pn}

\textbf{\citet{Blaetke2015}}: ``BioModel Engineering with Petri Nets''

\begin{itemize}
    \item Comprehensive survey of Bio-PN tools and applications
    \item Case studies: Signaling (MAPK), metabolism (TCA), gene circuits
    \item \textbf{Contribution}: Unified terminology, best practices
    \item \textbf{Gap identified}: ``No tool integrates metabolic kinetics with gene regulation''
\end{itemize}

\textbf{\citet{Liu2013}}: ``Multiscale Modeling with Hybrid Petri Nets''

\begin{itemize}
    \item Proposed framework for enzyme-catalyzed reactions + gene expression
    \item \textbf{Limitation}: Still requires two separate models (metabolism PN + regulation PN)
\end{itemize}

\textbf{\citet{Marwan2011}}: ``Petri Nets in Snoopy: A Unifying Framework''

\begin{itemize}
    \item Extended Snoopy with stochastic and hybrid PNs
    \item \textbf{Contribution}: Tool supporting multiple PN classes
    \item \textbf{Limitation}: No biochemical formula tracking, no weak independence analysis
\end{itemize}

\subsection{Gap Analysis}
\label{subsec:bio-pn-gap}

\textbf{What existing Bio-PNs provide}:

\begin{itemize}
    \item[$\checkmark$] Stoichiometric modeling (metabolic pathways)
    \item[$\checkmark$] Qualitative regulation (activation/inhibition arcs)
    \item[$\checkmark$] Hybrid dynamics (continuous + discrete)
    \item[$\checkmark$] Stochastic simulation (Gillespie algorithm)
\end{itemize}

\textbf{What is missing} (addressed by this thesis):

\begin{itemize}
    \item[$\times$] \textbf{Weak independence theory}: Parallel execution with shared catalysts/outputs
    \item[$\times$] \textbf{Heterogeneous coexistence}: Continuous + stochastic + timed + burst in single model
    \item[$\times$] \textbf{Arc-level regulatory semantics}: Hill equations on inhibitor arcs (not just constants)
    \item[$\times$] \textbf{Biochemical formula tracking}: Elemental composition, atomic conservation
    \item[$\times$] \textbf{Integrated metabolic-genetic models}: Single formalism spanning both layers
\end{itemize}

\section{Multi-Scale Biological Modeling}
\label{sec:multiscale-modeling}

\subsection{Genome-Scale Metabolic Models (GSMM)}
\label{subsec:gsmm}

\textbf{Flux Balance Analysis (FBA)} \citep{Varma1994}:

\begin{itemize}
    \item Models metabolism as \textbf{stoichiometric matrix} $S$ (metabolites $\times$ reactions)
    \item \textbf{Steady-state assumption}: $S \cdot v = 0$ (no accumulation)
    \item \textbf{Optimization}: Maximize biomass flux subject to constraints
    \item \textbf{Databases}: BiGG Models (7000+ curated GSMMs for bacteria, yeast, human)
\end{itemize}

\textbf{Advantages}:

\begin{itemize}
    \item \textbf{Scalability}: Models 1000+ reactions (entire metabolism)
    \item \textbf{Predictive power}: Predicts growth rates, gene essentiality
\end{itemize}

\textbf{Limitations}:

\begin{itemize}
    \item \textbf{No regulation}: Enzyme expression levels ignored (assumes all enzymes present)
    \item \textbf{No kinetics}: Flux directions but not rates (no time dynamics)
    \item \textbf{Steady-state only}: Cannot model transient responses
\end{itemize}

\subsection{Gene Regulatory Networks (GRN)}
\label{subsec:grn}

\textbf{Boolean networks} \citep{Kauffman1969}:

\begin{itemize}
    \item Genes are binary (on/off)
    \item Update rules: $\text{Gene}_i(t+1) = f(\text{Gene}_j(t), \text{Gene}_k(t), \ldots)$
    \item \textbf{Applications}: Cell differentiation, synthetic circuits
\end{itemize}

\textbf{Limitations}:

\begin{itemize}
    \item \textbf{No concentrations}: Cannot represent ``high vs. low ATP''
    \item \textbf{Discrete time}: Cannot model continuous dynamics
    \item \textbf{No metabolism}: Genes regulate each other, but metabolites absent
\end{itemize}

\textbf{Differential equation models}:

\begin{itemize}
    \item Genes/proteins as continuous variables
    \item Example: $\frac{d[X]}{dt} = k_{\text{synth}} - k_{\text{deg}} \cdot [X]$
    \item \textbf{Advantages}: Quantitative, captures kinetics
\end{itemize}

\textbf{Limitations}:

\begin{itemize}
    \item \textbf{Not compositional}: Adding genes requires rewriting equations
    \item \textbf{No visual topology}: Regulatory structure hidden in code
    \item \textbf{Separate from metabolism}: Metabolites not integrated
\end{itemize}

\subsection{Integrated Approaches}
\label{subsec:integrated-approaches}

\textbf{E-Cell} \citep{Tomita1999}:

\begin{itemize}
    \item Software platform simulating whole-cell models
    \item Combines multiple formalisms (ODEs, Gillespie, rule-based)
    \item \textbf{Contribution}: Demonstrated feasibility of multi-scale simulation
    \item \textbf{Limitation}: Not a formal modeling language (software tool, not formalism)
\end{itemize}

\textbf{Virtual Cell} \citep{Schaff1997}:

\begin{itemize}
    \item Spatial models with reaction-diffusion PDEs
    \item Compartments (cytoplasm, nucleus, membrane)
    \item \textbf{Contribution}: Spatial resolution
    \item \textbf{Limitation}: Computationally expensive, no formal analysis (P-invariants)
\end{itemize}

\textbf{Rule-based modeling} (BioNetGen, Kappa):

\begin{itemize}
    \item Molecular interactions as rewrite rules
    \item Compactly represents protein complexes with many states
    \item \textbf{Contribution}: Handles combinatorial complexity (phosphorylation sites)
    \item \textbf{Limitation}: Not Petri nets (different formalism), no elemental balance
\end{itemize}

\subsection{Why Current Approaches Fall Short}
\label{subsec:multiscale-shortcomings}

\textbf{Fundamental issue}: Existing multi-scale approaches either:

\begin{enumerate}
    \item \textbf{Separate models}: Build metabolism model + regulation model, integrate ad-hoc (no formal semantics)
    \item \textbf{Software-specific}: Tools provide integration (E-Cell, Virtual Cell) but no underlying formalism
    \item \textbf{Limited scope}: Focus on one layer (GSMM for metabolism, GRN for regulation)
\end{enumerate}

\textbf{This thesis provides}: A \textbf{formal Petri net extension} enabling integrated modeling with well-defined semantics, analyzability, and visual topology.

\section{Systems Biology Standards and Databases}
\label{sec:standards-databases}

\subsection{SBML (Systems Biology Markup Language)}
\label{subsec:sbml}

\textbf{\citet{Hucka2003}}: XML format for biochemical models

\begin{itemize}
    \item \textbf{Core}: Species, reactions, compartments, parameters
    \item \textbf{Extensions}: SBML-qual (qualitative GRNs), SBML-comp (model composition)
    \item \textbf{Support}: 300+ tools read/write SBML
\end{itemize}

\textbf{Advantages}:

\begin{itemize}
    \item \textbf{Interoperability}: Exchange models between tools
    \item \textbf{Standardization}: Widely adopted ($>10,000$ models)
\end{itemize}

\textbf{Limitations}:

\begin{itemize}
    \item \textbf{Not a formalism}: SBML is a data format, not a modeling language
    \item \textbf{Separate qualitative/quantitative}: SBML (kinetics) vs. SBML-qual (Boolean logic) are disconnected
    \item \textbf{No formal semantics}: Tool-dependent interpretation
\end{itemize}

\subsection{KEGG (Kyoto Encyclopedia of Genes and Genomes)}
\label{subsec:kegg}

\textbf{\citet{Kanehisa2000}}: Database of biological pathways

\begin{itemize}
    \item \textbf{KEGG COMPOUND}: 18,000+ metabolites (formulas, structures)
    \item \textbf{KEGG REACTION}: 11,000+ reactions (stoichiometry, EC numbers)
    \item \textbf{KEGG PATHWAY}: 500+ reference pathways (glycolysis, TCA, etc.)
    \item \textbf{REST API}: Programmatic access (Chapter~\ref{ch:kegg} uses this)
\end{itemize}

\textbf{Usage in modeling}:

\begin{itemize}
    \item Import pathways (reaction stoichiometry)
    \item Fetch compound formulas (\ce{C6H12O6} for glucose)
    \item Link to enzymes (EC numbers)
\end{itemize}

\textbf{Limitations}:

\begin{itemize}
    \item \textbf{Incomplete stoichiometry}: Often omits cofactors (\ce{H2O}, \ce{H+})
    \item \textbf{No kinetics}: Reaction topology only, no rate constants
\end{itemize}

\subsection{BRENDA (Enzyme Database)}
\label{subsec:brenda}

\textbf{\citet{Schomburg2004}}: Comprehensive enzyme kinetics

\begin{itemize}
    \item \textbf{2.7 million} parameter entries ($K_m$, $k_{cat}$, $K_i$, pH optima)
    \item \textbf{83,000+ enzymes} classified by EC number
    \item \textbf{Organism-specific}: Data from bacteria, yeast, human, etc.
    \item \textbf{SOAP API}: Programmatic access (Chapter~\ref{ch:brenda} uses this)
\end{itemize}

\textbf{Applications}:

\begin{itemize}
    \item Parameter inference for models ($K_m$ for glucose $\approx 0.1$ mM)
    \item Organism-specific parameterization (yeast vs. human kinetics)
\end{itemize}

\textbf{Challenges}:

\begin{itemize}
    \item \textbf{High variance}: Same enzyme, different labs $\to$ 10$\times$ differences
    \item \textbf{Incomplete coverage}: Not all enzymes have kinetic data
\end{itemize}

\subsection{Other Databases}
\label{subsec:other-databases}

\textbf{ChEBI} (Chemical Entities of Biological Interest):

\begin{itemize}
    \item Molecular structures, formulas, synonyms
    \item More detailed than KEGG (includes protonation states)
\end{itemize}

\textbf{SABIO-RK} (Reaction Kinetics Database):

\begin{itemize}
    \item Curated kinetic laws (full rate equations, not just $K_m$)
    \item Smaller than BRENDA (50,000 entries) but higher quality
\end{itemize}

\textbf{BiGG Models}:

\begin{itemize}
    \item Genome-scale metabolic reconstructions
    \item Flux constraints, gene-protein-reaction associations
\end{itemize}

\textbf{Integration in this thesis}:

\begin{itemize}
    \item \textbf{KEGG}: Formula retrieval (Chapter~\ref{ch:kegg})
    \item \textbf{BRENDA}: Parameter inference (Chapter~\ref{ch:brenda})
    \item \textbf{ChEBI}: Fallback for formulas (future work)
\end{itemize}

\section{Related Tools}
\label{sec:related-tools}

\subsection{Snoopy (Brandenburg University of Technology)}
\label{subsec:snoopy}

\textbf{\citet{Heiner2012}}: ``Snoopy - A Unifying Petri Net Tool''

\textbf{Supports}: Classical, timed, stochastic, hybrid Petri nets

\textbf{Features}:

\begin{itemize}
    \item Graphical editor (drag-and-drop)
    \item Simulation (ODE, Gillespie, hybrid)
    \item Model checking (CTL properties)
    \item SBML import/export
\end{itemize}

\textbf{Biological focus}: Signaling networks, metabolic pathways

\textbf{Strengths}:

\begin{itemize}
    \item Open-source, actively maintained
    \item Excellent visualization
    \item Structural analysis (P-invariants, reachability graph)
\end{itemize}

\textbf{Limitations}:

\begin{itemize}
    \item \textbf{No weak independence analysis}: Parallelism not exploited
    \item \textbf{No biochemical formulas}: Tokens are abstract
    \item \textbf{Limited regulation}: Inhibitor arcs have constant thresholds (no Hill equations)
    \item \textbf{Separate models for metabolism/regulation}: No unified formalism
\end{itemize}

\subsection{Cell Illustrator (University of Tokyo)}
\label{subsec:cell-illustrator}

\textbf{\citet{Nagasaki2010}}: Hybrid Petri net tool

\begin{itemize}
    \item \textbf{Hybrid dynamics}: Continuous + discrete + stochastic
    \item \textbf{Biological templates}: Gene expression, signaling cascades
    \item \textbf{3D visualization}: Network layout algorithms
\end{itemize}

\textbf{Strengths}:

\begin{itemize}
    \item Powerful hybrid simulation
    \item Large library of biological models
\end{itemize}

\textbf{Limitations}:

\begin{itemize}
    \item \textbf{Proprietary}: Commercial license required (not open-source)
    \item \textbf{Custom formalism}: Not standard Petri nets (tool-specific notation)
    \item \textbf{No KEGG/BRENDA integration}: Manual parameter entry
\end{itemize}

\subsection{COPASI (Complex Pathway Simulator)}
\label{subsec:copasi}

\textbf{\citet{Hoops2006}}: Biochemical system simulator

\begin{itemize}
    \item \textbf{ODE simulation}: Deterministic, stochastic (Gillespie), hybrid
    \item \textbf{Parameter estimation}: Fit models to experimental data
    \item \textbf{Optimization}: Metabolic flux optimization
    \item \textbf{SBML support}: Import/export
\end{itemize}

\textbf{Strengths}:

\begin{itemize}
    \item Widely used (5000+ citations)
    \item Parameter fitting (least-squares, genetic algorithms)
    \item Sensitivity analysis
\end{itemize}

\textbf{Limitations}:

\begin{itemize}
    \item \textbf{Not Petri nets}: Reaction network formalism (different from PNs)
    \item \textbf{No regulation arcs}: Inhibition encoded in rate laws (not topology)
    \item \textbf{Single-scale}: Metabolism-focused (gene regulation requires separate SBML-qual model)
\end{itemize}

\subsection{CellDesigner (Systems Biology Institute)}
\label{subsec:celldesigner}

\textbf{\citet{Funahashi2003}}: Visual SBML editor

\begin{itemize}
    \item \textbf{Graphical notation}: Process diagrams (SBGN standard)
    \item \textbf{SBML generation}: Automatically creates SBML from diagrams
    \item \textbf{Simulation}: Integrates with COPASI, CellML
\end{itemize}

\textbf{Strengths}:

\begin{itemize}
    \item User-friendly visual editor
    \item SBML standardization
\end{itemize}

\textbf{Limitations}:

\begin{itemize}
    \item \textbf{Not Petri nets}: Process diagrams (different semantics)
    \item \textbf{No formal analysis}: Cannot compute P-invariants, reachability
    \item \textbf{Separate regulation}: Metabolic and regulatory diagrams disconnected
\end{itemize}

\subsection{Charlie (Model Checker)}
\label{subsec:charlie}

\textbf{\citet{Heiner2015}}: Verification tool for Petri nets

\begin{itemize}
    \item \textbf{Model checking}: CTL, LTL temporal logic properties
    \item \textbf{Symbolic methods}: BDDs for state space exploration
    \item \textbf{Biological properties}: Verifies liveness, boundedness, reachability
\end{itemize}

\textbf{Strengths}:

\begin{itemize}
    \item Formal verification (prove properties hold)
    \item Handles large state spaces (symbolic encoding)
\end{itemize}

\textbf{Limitations}:

\begin{itemize}
    \item \textbf{Verification only}: Not a modeling tool (requires Snoopy for model creation)
    \item \textbf{No kinetics}: Qualitative properties only (not quantitative simulation)
\end{itemize}

\subsection{Comparison Summary}
\label{subsec:tools-comparison}

Table~\ref{tab:tools-comparison} compares related tools with SHYpn (this thesis).

\begin{table}[htbp]
\centering
\caption{Comparison of biological modeling tools.}
\label{tab:tools-comparison}
\small
\begin{tabular}{lcccccc}
\toprule
\textbf{Tool} & \textbf{Type} & \textbf{Hybrid} & \textbf{Regulation} & \textbf{KEGG/} & \textbf{Weak} & \textbf{Formula} \\
 &  &  &  & \textbf{BRENDA} & \textbf{Indep.} & \textbf{Track.} \\
\midrule
Snoopy & Petri net & Yes & Limited & No & No & No \\
Cell Ill. & Hybrid PN & Yes & Yes & No & No & No \\
COPASI & ODE/Gill. & Yes & Rate laws & No & N/A & No \\
CellDesigner & SBML ed. & Via tools & Proc. diag. & No & N/A & No \\
Charlie & Model check & Yes & Yes & No & No & No \\
\midrule
\textbf{SHYpn} & Ext. Bio-PN & Yes & Arc-level & Yes & Yes & Yes \\
\bottomrule
\end{tabular}
\end{table}

\textbf{Key differentiator}: SHYpn is the \textbf{first tool} providing:

\begin{enumerate}
    \item Weak independence-based parallelism
    \item Arc-level regulation (Hill equations on inhibitor arcs)
    \item Biochemical formula tracking (elemental balance)
    \item Automated KEGG/BRENDA integration
\end{enumerate}

\section{Theoretical Foundations}
\label{sec:theoretical-foundations}

\subsection{Concurrency Theory}
\label{subsec:concurrency-theory}

\textbf{Independence of transitions} (classical definition):

Transitions $t_1, t_2$ are \textbf{independent} if:

\begin{equation}
(\bullet t_1 \cup t_1\bullet) \cap (\bullet t_2 \cup t_2\bullet) = \emptyset
\end{equation}

(No shared places in neighborhoods)

\textbf{Consequence}: Can fire \textbf{concurrently} without conflict

\textbf{Diamond property} \citep{Mazurkiewicz1987}:

If $t_1, t_2$ are independent and both enabled at $M$:

\begin{align}
M &\xrightarrow{t_1} M_1 \xrightarrow{t_2} M_{12} \\
M &\xrightarrow{t_2} M_2 \xrightarrow{t_1} M_{12}
\end{align}

(Same final state $M_{12}$ regardless of order)

\textbf{Limitation for biology}:

\begin{itemize}
    \item Enzymes catalyze multiple reactions (shared place) $\to$ Not independent by classical definition
    \item But \textbf{biologically}, these reactions don't conflict (enzyme is preserved)
\end{itemize}

\textbf{This thesis contribution}: Weak independence relaxes the condition to allow shared test arcs and outputs while preserving diamond property.

\subsection{Process Calculi}
\label{subsec:process-calculi}

\textbf{CCS} (Calculus of Communicating Systems) \citep{Milner1980}:

\begin{itemize}
    \item Algebraic formalism for concurrent processes
    \item Composition operator ($P | Q$) runs $P$, $Q$ in parallel
\end{itemize}

\textbf{$\pi$-calculus} \citep{Milner1992}:

\begin{itemize}
    \item Extends CCS with mobile processes (channels can be passed)
\end{itemize}

\textbf{Stochastic $\pi$-calculus} \citep{Priami1995}:

\begin{itemize}
    \item Adds stochastic rates to reactions
    \item Used in computational biology (BioAmbients, Beta-binders)
\end{itemize}

\textbf{Relation to Petri nets}:

\begin{itemize}
    \item Process calculi and Petri nets are \textbf{equivalent in expressiveness} (both Turing-complete)
    \item Petri nets are \textbf{graphical} (visual topology), process calculi are \textbf{textual} (algebraic)
    \item Choice: Petri nets preferred in systems biology for \textbf{visual reasoning}
\end{itemize}

\subsection{Chemical Reaction Network Theory}
\label{subsec:crn-theory}

\textbf{\citet{Feinberg1987}}: Deficiency theory

\begin{itemize}
    \item Analyzes steady states of reaction networks
    \item \textbf{Deficiency zero theorem}: Networks with deficiency $\delta=0$ have unique positive steady state
\end{itemize}

\textbf{Relationship to Petri nets}:

\begin{itemize}
    \item Reaction networks are \textbf{continuous Petri nets} with mass action kinetics
    \item Stoichiometric matrix $S$ corresponds to PN incidence matrix
\end{itemize}

\textbf{Limitation}:

\begin{itemize}
    \item Mass action only (no Michaelis-Menten, Hill)
    \item No gene regulation (reactions among metabolites only)
\end{itemize}

\subsection{Stochastic Simulation}
\label{subsec:stochastic-simulation}

\textbf{\citet{Gillespie1977}}: Stochastic Simulation Algorithm (SSA)

\begin{itemize}
    \item \textbf{Exact algorithm} for simulating chemical master equation (CME)
    \item Two random numbers per step:
    \begin{enumerate}
        \item When does next reaction occur? ($\tau \sim$ Exponential)
        \item Which reaction fires? (Weighted choice by propensity)
    \end{enumerate}
\end{itemize}

\textbf{$\tau$-leaping} \citep{Gillespie2001}:

\begin{itemize}
    \item Approximate method: Fire multiple reactions per step
    \item Faster but loses exactness
\end{itemize}

\textbf{Hybrid methods} \citep{Haseltine2002}:

\begin{itemize}
    \item Fast reactions: Continuous (ODE)
    \item Slow reactions: Stochastic (Gillespie)
    \item \textbf{Challenge}: Partitioning (which reactions are fast?)
\end{itemize}

\textbf{This thesis contribution}: Hybrid scheduler coordinates ODE + Gillespie + Timed + Burst engines without manual partitioning (user specifies transition type explicitly).

\section{Gap Summary}
\label{sec:gap-summary}

\textbf{Existing work provides}:

\begin{itemize}
    \item[$\checkmark$] Petri net theory (classical, timed, stochastic, hybrid)
    \item[$\checkmark$] Biological applications (metabolic pathways, signaling networks)
    \item[$\checkmark$] Tools (Snoopy, COPASI, CellDesigner)
    \item[$\checkmark$] Databases (KEGG, BRENDA, SBML)
    \item[$\checkmark$] Simulation algorithms (ODE, Gillespie, hybrid)
\end{itemize}

\textbf{What is missing} (this thesis addresses):

\begin{enumerate}
    \item \textbf{Weak independence theory}: Formal framework for parallel execution with shared catalysts/outputs
    \item \textbf{Heterogeneous transition coexistence}: Continuous + stochastic + timed + burst in single model with unified semantics
    \item \textbf{Arc-level regulatory semantics}: Hill equations, threshold formulas on arcs (not external code)
    \item \textbf{Biochemical formula tracking}: Elemental composition, atomic conservation verification
    \item \textbf{Integrated metabolic-genetic formalism}: Single framework spanning biochemistry and gene regulation
    \item \textbf{Tool integration}: Automated KEGG (formulas) + BRENDA (parameters) enrichment
\end{enumerate}

\textbf{Positioning}: This thesis is the \textbf{first} to provide all six capabilities in a unified Petri net formalism.

\section{Summary}
\label{sec:ch2-summary}

\textbf{This chapter surveyed foundational concepts and related work}:

\begin{itemize}
    \item \textbf{Section~\ref{sec:pn-fundamentals}}: Classical Petri nets (5-tuple, enabling, firing, properties)
    \item \textbf{Section~\ref{sec:pn-extensions}}: Extensions (colored, timed, stochastic, hybrid)
    \item \textbf{Section~\ref{sec:bio-pn}}: Biological Petri nets (Reddy 1993 $\to$ present, Snoopy tool)
    \item \textbf{Section~\ref{sec:multiscale-modeling}}: Multi-scale modeling (GSMM, GRN, E-Cell, Virtual Cell)
    \item \textbf{Section~\ref{sec:standards-databases}}: Standards and databases (SBML, KEGG, BRENDA)
    \item \textbf{Section~\ref{sec:related-tools}}: Tools comparison (Snoopy, COPASI, CellDesigner, Charlie)
    \item \textbf{Section~\ref{sec:theoretical-foundations}}: Theoretical foundations (concurrency, process calculi, Gillespie)
    \item \textbf{Section~\ref{sec:gap-summary}}: Gap analysis (what is missing)
\end{itemize}

\textbf{Key findings}:

\begin{itemize}
    \item Petri nets are \textbf{well-suited} for biological modeling (visual, compositional, analyzable)
    \item Existing Bio-PNs model \textbf{metabolism or regulation, not both} in integrated way
    \item \textbf{No formalism} provides weak independence, arc-level regulation, and formula tracking together
    \item \textbf{Tools} (Snoopy, COPASI) are powerful but lack these features
\end{itemize}

\textbf{Next chapter} (Chapter~\ref{ch:integration-challenge}) presents the \textbf{integration challenge} in depth, deriving formal requirements from the lac operon system.

\textbf{Transition}: Background established $\to$ Now define the problem precisely.
